\documentclass[12pt]{article}

\usepackage{newtxtext, amsthm, amssymb, amsmath, newtxmath}
\usepackage[margin=1in]{geometry}
\usepackage{setspace}
\doublespace
\title{\textbf{AMS 572 Project Report}}
\author{Benjamin Minkin, Marcus Minutello, Jacob White}
\date{Due: December 5, 2025 \@ 10:00am}
\begin{document} 
	\maketitle
	\section{Introduction}
	Nearly two decades since its introduction in 2009, Bitcoin has gained mainstream traction as an
	investment asset. This digital currency is held by major heads funds[1], Fortune 500 companies[2],
	and national banks[3]. Additionally, Bitcoin has gained acceptance as an everyday currency with many
	retailers now accepting it as payment. Most incredibly, Bitcoin was even declared a legal-tender currency
	of El Salvador in September 2021[4]. As Bitcoin has become more accepted, retail investors calmored to buy 
	and trade this new asset. Today, you can hold Bitcoin in digital wallets, investment and retirement accounts[5],
	as well as purchasing shares of Bitcoin-focused exchange-traded funds (ETFs)[6].
	
	Moreover, the meteoric rise in price per bitcoin has spawned a new class of digital assets called cryptocurrencies. 
	Aside from Bitcoin, these cryptocurrencies follow one of two paradigms. The first paradigm describes coins that focus 
	on improving some aspects of the technology such as energy efficiency, transaction speed, or added utility. These coins 
	utilize their own blockchain networking structure leading to their name as “native coins”. The second type of coin attempts 
	to capitalize on the speculation and excitement around Bitcoin’s price by creating eye-catching imitations. These coins use
	existing blockchain networks and are often themed around a specific internet trend or personality to garner attention. In 
	online circles, these types of coins have earned the title of “meme coin” as they make no novel improvements to the 
	underlying structure.  
	
	Due to its mainstream adoption in traditional financial markets by both institutional and retail investors, does Bitcoin’s 
	price follow stock market movements? Is Bitcoin truly a new investment asset, or is its price so deeply entangled with 
	traditional investments that it has lost its novelty? The answer to this question has serious ramifications for portfolio 
	allocations. If Bitcoin’s movements are unrelated to the stock market, it can be a valuable tool to diversify portfolios 
	and reduce risk. However, if Bitcoin’s price follows the movements of the stock market, its utility as an investment asset 
	mirrors similar stocks in the digital sphere. 

	How do the price movements of Bitcoin proportionally affect Native or Meme coins? Are investors more excited about the underlying 
	technology powering digital assets or are they seeking an avenue to wildly speculate? The answer to this question will help unravel 
	the motivation fueling the cryptocurrency boom. This is because of the dual purpose of currency, being a store-of-value and a mode 
	of facilitating transactions. Both of these purposes benefit from the price being predictable and stable. If the meme coins are more 
	sensitive to changes in Bitcoin, it highlights their use as a speculative asset. Likewise, this result highlights the utility of native
	coins and investors bet on the underlying technological improvements. 

	\subsection{Questions of Interest}
	
	Due to its mainstream adoption in traditional financial markets by both institutional and retail investors, does Bitcoin’s price follow stock market
	movements? Is Bitcoin truly a new investment asset, or is its price so deeply entangled with traditional investments that it has lost its novelty? 
	The answer to this question has serious ramifications for portfolio allocations. If Bitcoin’s movements are unrelated to the stock market, it can 
	be a valuable tool to diversify portfolios and reduce risk. However, if Bitcoin’s price follows the movements of the stock market, its utility as 
	an investment asset mirrors similar stocks in the digital sphere. 
	
	How do the price movements of Bitcoin proportionally affect Native or Meme coins? Are investors more excited about the underlying technology powering
	digital assets or are they seeking an avenue to wildly speculate? The answer to this question will help unravel the motivation fueling the
	cryptocurrency boom. This is because of the dual purpose of currency, being a store-of-value and a mode of facilitating transactions. 
	Both of these purposes benefit from the price being predictable and stable. If the meme coins are more sensitive to changes in Bitcoin, 
	it highlights their use as a speculative asset. Likewise, this result highlights the utility of native coins and investors bet on the 
	underlying technological improvements. 
	
	\subsection{Description of the Data}
	
	The data for this analysis was imported from Yahoo Finance using the quantmod r package. This data contains the date, daily opening price, daily high 
	price, daily closing price, and daily volume for each asset specified. The range of dates collected for this analysis ranges from November 1, 2023
	to November 18, 2025. This data was then transformed to isolate the log-daily returns over this time period. Additionally, a dummy variable was
	added to categorize between the native coins and the meme coins. When combining the SPY and BTC data, non-trading days such as holidays and
	weekends were removed. 
	
	\section{Background}
	In this section, we review some basic mathematical machinery with regard to financial time series. 
	\subsection{Introduction to Financial Time Series}
	By a time series we mean a sequence of observations taken at regular time intervals. More formally, a \textbf{time series}  is simply a sequence $\{x_t\}_{t \in T}$ of real numbers $x_t$, for $T = \{t_0, t_1, t_2, t_3, \dots, t_n\}$ some finite indexing set of time increments over a certain period\footnote{In the case of real-world data, this indexing set is finite. However, generally we can make $T$ to be any finite, countably infinite, or even uncountably infinite subset of $\mathbb{R}$ (the real numbers).}. In the case of finance, our sequence $\{x_t\}_{t \in T}$ (which we may write equivalently as $x_t$ for brevity, where the time index $T$ is implicitly understood) represents the evolution of the price of a particular asset over time.  To our notation consistent, we denote the price of a particular asset at the end of a trading day $t$ by $P_t$ . 
	
	
	When we study the evolution of an asset price, it is often useful to measure changes in price by computing the \textbf{daily log return} time series $r_t$ of $P_t$, which is given by the formula 
		\begin{equation}
			\label{eq:logret}
			r_t = \log\left(\frac{P_t}{P_{t - 1}}\right)
		\end{equation} 
	\subsection{Correlation and Time Series}
	\subsection{Linear Regression on Financial Time Series}
	\subsection{Autocorrelation, Heteroskedasticity, and Robust Inference}
	
	\section{Question 1}
	\section{Question 2}

	\section{Conclusion}
	
	\bibliographystyle{plain} 
	\bibliography{references}
\end{document}